\chapter{Osciladores mecânicos e a medida do tempo}\label{ch:g12:oscillators-and-time}

Um relógio de pêndulo faz tique: cada tique é um pêndulo de latão cruzando a vertical,
repetido quinze milhões de vezes desde o Ano-Novo. Uma lâmina de quartzo no seu relógio
de pulso verga \num{32768} vezes por segundo; um \omterm{def:g11:fundamental-interactions:ladder}{átomo} de césio zumbe nove bilhões de
vezes mais depressa. Tudo que balança pode marcar o tempo: este capítulo acompanha o
balanço, de Huygens aos \omterm{def:g10:universal-gravitation:satellite}{satélites}.

\section{Oscilações livres}

\begin{definition}[Oscilador, período, frequência, amplitude]\label{def:g12:oscillators-and-time:oscillator}
Um \emph{oscilador mecânico}\index{oscilador} é um sistema que, afastado de uma
\emph{posição de equilíbrio}\index{posição de equilíbrio} estável e solto, vai e volta
em torno dela: são as \emph{oscilações livres}\index{oscilações livres}. O movimento
se repete depois do \emph{período}\index{período} $T_0$ (\unit{s}); a
\emph{frequência}\index{frequência} é $f_0 = 1/T_0$ (\unit{Hz},
\cref{ch:g12:mechanical-waves}); a \emph{amplitude}\index{amplitude} é o afastamento
máximo.
\end{definition}

\begin{example}[Um zoológico de osciladores]\label{ex:g12:oscillators-and-time:zoo}
Um balanço de parquinho: $T_0 \approx \qty{3}{s}$. A corda lá de um violão:
$f_0 = \qty{110}{Hz}$. Um cristal de quartzo: \qty{32768}{Hz}. A oscilação do césio que
define o segundo: cerca de \qty{9.2e9}{Hz}. Dez \omterm{def:g10:orders-of-magnitude:oom}{ordens de grandeza}, um único ofício:
contar.
\end{example}

\section{O pêndulo simples}

\begin{definition}[Pêndulo simples]\label{def:g12:oscillators-and-time:pendulum}
Um \emph{pêndulo simples}\index{pêndulo simples}: um pequeno corpo de massa $m$ preso a
um fio inextensível de comprimento $L$ e massa desprezível, oscilando num plano
vertical. Das duas \omterm{def:g10:inertia:force}{forças} sobre o corpo (\cref{ch:g12:newtons-laws}), a tensão é
perpendicular ao movimento, enquanto a componente do \omterm{def:g10:universal-gravitation:weight}{peso} ao longo da \omterm{def:g10:relative-motion:trajectory}{trajetória} aponta
de volta para o ponto mais baixo: o \omterm{def:g10:universal-gravitation:weight}{peso} é a \emph{\omterm{def:g10:inertia:force}{força} restauradora}%
\index{força restauradora}, que se opõe ao afastamento.
\end{definition}

\begin{omfigure}
\begin{tikzpicture}[scale=0.8, font=\small]
  \coordinate (O) at (0,0); \draw (-0.6,0) -- (0.6,0);
  \foreach \x in {-0.5,-0.25,0,0.25,0.5} \draw (\x,0) -- ++(0.14,0.14);
  \fill (O) circle (1.5pt); \draw[dashed] (O) -- (0,-3.4);
  \draw[dashed, omExo] (-120:3.0) arc (-120:-60:3.0);
  \draw (O) -- (-60:3.0) node[midway, above right] {$L$};
  \draw (-90:1.0) arc (-90:-60:1.0); \node at (-75:1.32) {$\theta$};
  \coordinate (B) at (1.5,-2.598); \fill[omDef] (B) circle (2.5pt);
  \draw[-Stealth, omDef, thick] (B) -- ++(-0.75,1.299) node[above left] {$\vect T$};
  \draw[-Stealth, thick] (B) -- ++(0,-1.8) node[below] {$\vect P$};
  \draw[-Stealth, omThm, thick] (B) -- ++(-0.779,-0.45) node[left] {$\vect{P_t}$};
  \draw[dashed, omExo] (0.721,-3.048) -- (1.5,-4.398)
    (B) -- ++(0.779,-1.35) (2.279,-3.948) -- (1.5,-4.398);
\end{tikzpicture}

{\small \omterm{def:g10:inertia:force}{Forças} sobre o corpo: a tensão $\vect T$ é perpendicular ao movimento; a parte
tangencial $\vect{P_t}$ do \omterm{def:g10:universal-gravitation:weight}{peso} sempre aponta de volta para a vertical --- a \omterm{def:g12:oscillators-and-time:pendulum}{força
restauradora}.}
\end{omfigure}

\begin{proposition}[Período do pêndulo simples]\label{prop:g12:oscillators-and-time:pendulum-period}
\index{pêndulo simples!período}
Para \omterm{def:g12:oscillators-and-time:oscillator}{amplitudes} pequenas (abaixo de uns $15^\circ$), o \omterm{def:g12:oscillators-and-time:oscillator}{período} é
\[
T_0 = 2\pi \sqrt{\frac{L}{g}}, \qquad g = \qty{9.81}{m/s^2}.
\]
Ele não depende nem da massa do corpo nem --- enquanto o ângulo permanecer pequeno ---
da \omterm{def:g12:oscillators-and-time:oscillator}{amplitude}: as oscilações pequenas são \emph{isócronas}\index{isocronismo}, todas
levando o mesmo tempo.
\end{proposition}
\admitted

\begin{remark}[Análise dimensional, e onde mora o $2\pi$]\label{rem:g12:oscillators-and-time:dimensional}
A fórmula pode ser adivinhada. A partir de $m$ (\unit{kg}), $L$ (\unit{m}) e $g$
(\unit{m/s^2}), a única combinação com \omterm{def:g10:orders-of-magnitude:unit}{unidade} de tempo é $\sqrt{L/g}$ --- e a massa
\emph{não pode} entrar, já que nenhum outro dado carrega um \omterm{def:g10:orders-of-magnitude:unit}{quilograma} para cancelá-la.
As dimensões não conseguem dar o número puro da frente: o $2\pi$ vem de resolver a
equação do movimento, o que se faz no volume do primeiro ano de graduação. E a condição
de ângulo pequeno também não é opcional: a $30^\circ$ o \omterm{def:g12:oscillators-and-time:oscillator}{período} verdadeiro é $2\%$
maior, e o isocronismo falha.
\end{remark}

\begin{example}[O pêndulo de segundos]\label{ex:g12:oscillators-and-time:seconds}
Que comprimento bate o segundo, $T_0 = \qty{2.00}{s}$, um tique para cada lado?
Resolvendo, $L = g\,T_0^2/4\pi^2 = 9.81 \times 2.00^2/4\pi^2 \approx \qty{0.994}{m}$.
Quase um \omterm{def:g10:orders-of-magnitude:unit}{metro}: os relógios de pêndulo são altos porque os segundos são longos.
\end{example}

\section{O oscilador massa--mola}

\begin{definition}[Força restauradora de uma mola]\label{def:g12:oscillators-and-time:spring}
Um carrinho de massa $m$ desliza sobre um trilho horizontal sem \omterm{def:g10:inertia:inventory}{atrito}, preso a uma
mola de \emph{constante elástica}\index{constante elástica} $k$ (\unit{N/m}). Seja $x$
o afastamento dele do equilíbrio. Esticada ou comprimida, a mola puxa ou empurra o
carrinho de volta com a \omterm{def:g10:inertia:force}{força} $F = -kx$: proporcional ao afastamento e oposta a ele ---
de novo uma \omterm{def:g12:oscillators-and-time:pendulum}{força restauradora}, agora fornecida pela elasticidade.
\end{definition}

\begin{proposition}[Equação do movimento e a solução dela]\label{prop:g12:oscillators-and-time:spring-period}
\index{oscilador massa--mola}
A \omterm{thm:g12:newtons-laws:second}{segunda lei de Newton} ao longo do trilho (\cref{ch:g12:newtons-laws}) dá
$m \frac{d^2x}{dt^2} = -kx$, ou seja, $\frac{d^2x}{dt^2} = -\frac{k}{m}\,x$. Sejam
quais forem as constantes $A$ e $\varphi$,
\[
x(t) = A \cos\!\left(\frac{2\pi t}{T_0} + \varphi\right),
\qquad T_0 = 2\pi\sqrt{\frac{m}{k}},
\]
satisfaz essa equação: oscilação de \omterm{def:g12:oscillators-and-time:oscillator}{período} $T_0$, independente da \omterm{def:g12:oscillators-and-time:oscillator}{amplitude}.
\end{proposition}

\begin{proof}
Derivando, $\frac{dx}{dt} = -\frac{2\pi}{T_0}A \sin(\frac{2\pi t}{T_0} + \varphi)$; de
novo, $\frac{d^2x}{dt^2} = -(\frac{2\pi}{T_0})^2 A\cos(\frac{2\pi t}{T_0} + \varphi) =
-(\frac{2\pi}{T_0})^2\,x(t)$, que é $-(k/m)\,x$ exatamente quando
$(2\pi/T_0)^2 = k/m$.
\end{proof}

\begin{remark}[Todos os movimentos que existem]\label{rem:g12:oscillators-and-time:uniqueness}
Que \emph{todo} movimento do carrinho é uma dessas funções fica admitido aqui e é
demonstrado no volume do primeiro ano de graduação. E não há $g$ em $2\pi\sqrt{m/k}$:
mais rígida oscila mais rápido, mais pesada mais devagar, e o mesmo vale na Lua.
\end{remark}

\begin{definition}[Amplitude e fase]\label{def:g12:oscillators-and-time:phase}
Em $x(t) = A\cos(2\pi t/T_0 + \varphi)$ a constante positiva $A$ é a \omterm{def:g12:oscillators-and-time:oscillator}{amplitude} --- os
extremos são $x = \pm A$ --- e $2\pi t/T_0 + \varphi$ é a \emph{fase}\index{fase}, que
cresce de $2\pi$ a cada \omterm{def:g12:oscillators-and-time:oscillator}{período}; a \emph{fase inicial}\index{fase inicial} $\varphi$
registra a partida: solto em repouso de $x = A$, tem-se $\varphi = 0$. Por exemplo,
$m = \qty{0.20}{kg}$ numa mola de $k = \qty{20}{N/m}$ dá
$T_0 = 2\pi\sqrt{0.20/20} = \qty{0.63}{s}$; puxado até \qty{5.0}{cm} e solto,
$x(t) = 5.0\cos(2\pi t/0.63)$ em centímetros.
\end{definition}

\section{Energia, amortecimento, ressonância}

\begin{proposition}[Energia do oscilador]\label{prop:g12:oscillators-and-time:energy}
Uma mola esticada de $x$ armazena a \omterm{def:g10:energy-conservation:energy}{energia} potencial elástica $E_p = \tfrac12 kx^2$
(justificada no \cref{ch:g12:work-and-energy}). Durante uma oscilação sem \omterm{def:g10:inertia:inventory}{atrito}, a
\omterm{def:g11:mechanical-energy:em}{energia mecânica}
\[
E_m = E_k + E_p = \tfrac12 m v^2 + \tfrac12 k x^2 = \tfrac12 k A^2
\]
é constante: toda elástica nos extremos, toda cinética no equilíbrio. O pêndulo toca o
mesmo dueto entre $E_k$ e a $E_p$ gravitacional (\cref{ch:g11:mechanical-energy}).
\end{proposition}

\begin{proof}
Com $v = \frac{dx}{dt}$, a derivada de $E_m$ é
$mv\frac{dv}{dt} + kx\frac{dx}{dt} = v\,(m\frac{d^2x}{dt^2} + kx) = 0$ pela equação do
movimento; o valor $\tfrac12 kA^2$ se lê num extremo.
\end{proof}

\begin{omfigure}
\begin{tikzpicture}
\begin{axis}[omaxis, width=8.6cm, height=4.6cm, xlabel={$t/T_0$},
    ylabel={energia $/\,E_m$}, xmin=0, xmax=1, ymin=0, ymax=1.3,
    legend pos=south east]
  \addplot[omDef, thick, domain=0:1, samples=120] {cos(360*x)^2};
  \addplot[omProp, thick, domain=0:1, samples=120] {sin(360*x)^2};
  \addplot[omThm, dashed, thick, domain=0:1, samples=2] {1};
  \legend{$E_p$, $E_k$, $E_m$}
\end{axis}
\end{tikzpicture}

{\small Um \omterm{def:g12:oscillators-and-time:oscillator}{período}, com partida em $x = A$: a \omterm{def:g10:energy-conservation:energy}{energia} vai e volta duas vezes da forma
elástica para a cinética, com a soma cravada em $\tfrac12 kA^2$.}
\end{omfigure}

\begin{definition}[Oscilações amortecidas]\label{def:g12:oscillators-and-time:damping}
O \omterm{def:g10:inertia:inventory}{atrito} drena \omterm{def:g11:mechanical-energy:em}{energia mecânica} de todo \omterm{def:g12:oscillators-and-time:oscillator}{oscilador} real: as \omterm{def:g12:oscillators-and-time:oscillator}{oscilações livres} dele são
\emph{amortecidas}\index{amortecimento}. Com amortecimento fraco o sistema ainda
oscila, com uma \omterm{def:g12:oscillators-and-time:oscillator}{amplitude} que morre mas um tempo de repetição quase inalterado, o
\emph{pseudoperíodo}\index{pseudoperíodo}, próximo de $T_0$; com amortecimento forte
(mel) ele volta rastejando ao equilíbrio sem nunca passar do ponto; a \omterm{def:g10:energy-conservation:energy}{energia},
proporcional ao quadrado da \omterm{def:g12:oscillators-and-time:oscillator}{amplitude}, decai virando calor.
\end{definition}

\begin{omfigure}
\begin{tikzpicture}
\begin{axis}[omaxis, width=8.6cm, height=4.6cm, xlabel={$t$ (\unit{s})},
    ylabel={$x$ (\unit{cm})}, xmin=0, xmax=10, ymin=-2.6, ymax=2.6,
    legend pos=south east]
  \addplot[omExo, thin, domain=0:10, samples=240] {2*cos(deg(pi*x))};
  \addplot[omDef, thick, domain=0:10, samples=240] {2*exp(-x/4)*cos(deg(pi*x))};
  \legend{sem amortecimento, amortecida}
  \addplot[omThm, dashed, domain=0:10, samples=60, forget plot] {2*exp(-x/4)};
  \addplot[omThm, dashed, domain=0:10, samples=60, forget plot] {-2*exp(-x/4)};
\end{axis}
\end{tikzpicture}

{\small Oscilações sem \omterm{def:g12:oscillators-and-time:damping}{amortecimento} e com \omterm{def:g12:oscillators-and-time:damping}{amortecimento} fraco: quase o mesmo \omterm{def:g12:oscillators-and-time:oscillator}{período}
(\qty{2.0}{s} aqui), com a \omterm{def:g12:oscillators-and-time:oscillator}{amplitude} amortecida morrendo dentro de um envelope
exponencial que encolhe (tracejado).}
\end{omfigure}

\begin{definition}[Oscilações forçadas e ressonância]\label{def:g12:oscillators-and-time:resonance}
Empurre um \omterm{def:g12:oscillators-and-time:oscillator}{oscilador} periodicamente --- \emph{oscilações forçadas}%
\index{oscilações forçadas} --- e, depois de um breve transitório, ele oscila na
\omterm{def:g12:oscillators-and-time:oscillator}{frequência} $f$ \emph{do excitador}, e não na dele. A \omterm{def:g12:oscillators-and-time:oscillator}{amplitude} dele é máxima quando $f$
se aproxima da \emph{frequência própria}\index{frequência própria} $f_0$ do \omterm{def:g12:oscillators-and-time:oscillator}{oscilador}
livre: é a \emph{ressonância}\index{ressonância}, com cada empurrão chegando em
compasso com o movimento e injetando \omterm{def:g10:energy-conservation:energy}{energia} a cada ciclo; quanto mais fraco o
\omterm{def:g12:oscillators-and-time:damping}{amortecimento}, mais alto e mais estreito o pico.
\end{definition}

\begin{omfigure}
\begin{tikzpicture}
\begin{axis}[omaxis, width=8.6cm, height=4.6cm, xlabel={$f/f_0$},
    ylabel={amplitude (u.a.)}, xmin=0, xmax=2.2, ymin=0, ymax=5.6,
    legend pos=north east]
  \addplot[omThm, thick, domain=0:2.2, samples=200] {1/sqrt((1-x^2)^2 + (x/5)^2)};
  \addplot[omDef, thick, domain=0:2.2, samples=120] {1/sqrt((1-x^2)^2 + (x/1.2)^2)};
  \legend{amortecimento fraco, amortecimento forte}
  \addplot[omExo, dashed, forget plot] coordinates {(1,0) (1,5.4)};
\end{axis}
\end{tikzpicture}

{\small Curva de \omterm{def:g12:oscillators-and-time:resonance}{ressonância}: \omterm{def:g12:oscillators-and-time:oscillator}{amplitude} forçada contra a \omterm{def:g12:oscillators-and-time:oscillator}{frequência} do excitador. Perto
de $f = f_0$ a resposta se ergue como uma torre, tanto mais alta quanto mais fraco o
\omterm{def:g12:oscillators-and-time:damping}{amortecimento}.}
\end{omfigure}

\begin{example}[Ressonância para o bem e para o mal]\label{ex:g12:oscillators-and-time:resonance-zoo}
Uma mãe empurrando um balanço empurra na \omterm{def:g12:oscillators-and-time:resonance}{frequência própria} dele --- \omterm{def:g12:oscillators-and-time:resonance}{ressonância} a
serviço. Em \num{2000} a recém-inaugurada Millennium Bridge, em Londres, balançou de
modo alarmante quando os passos da multidão se travaram numa \omterm{def:g12:oscillators-and-time:resonance}{frequência própria} perto
de \qty{1}{Hz}; ela ficou dois anos fechada para receber amortecedores. Todo relógio de
quartzo mantém o cristal dele cantando ao excitá-lo em \omterm{def:g12:oscillators-and-time:resonance}{ressonância}, a \qty{32768}{Hz}.
\end{example}

\section{Medir o tempo}

\begin{method}[Cronometrar um período]\label{met:g12:oscillators-and-time:timing}
Nunca cronometre uma única oscilação: dispare o cronômetro quando o corpo cruzar a
vertical (o ponto mais rápido, o mais fácil de julgar), conte $N = 50$ oscilações e
divida por $N$ --- um erro de reação de \qty{0.2}{s} vira \qty{4}{ms} em $T_0$.
Repita, e compare as tomadas.
\end{method}

\begin{example}[Três séculos de tique]\label{ex:g12:oscillators-and-time:clocks}
Huygens construiu o primeiro relógio de pêndulo em \num{1657}: o isocronismo levou os
relógios de um quarto de hora de desvio por dia para segundos, e os pêndulos
compensados, para frações de segundo por dia. O relógio de pulso de quartzo
(\num{1969}) trocou a haste por um cristal que verga a \qty{32768}{Hz}: décimos de
segundo por mês. Desde \num{1967} o próprio segundo é \emph{definido} por uma oscilação
atômica: a duração de \num{9192631770} \omterm{def:g12:oscillators-and-time:oscillator}{períodos} da radiação de \omterm{def:g10:signals-and-waves:wave}{micro-ondas} de uma
transição interna do césio-133 (\cref{ch:g12:quantum-world}). Os melhores relógios de
césio erram um segundo em cem milhões de anos.
\end{example}

\begin{remark}[Por que caçar o nanossegundo]\label{rem:g12:oscillators-and-time:gps}
Um receptor de GPS acha a posição dele cronometrando sinais de rádio vindos de
\omterm{def:g10:universal-gravitation:satellite}{satélites}; a luz percorre \qty{30}{cm} por nanossegundo, de modo que um erro de relógio
de um microssegundo o desloca em \qty{300}{m}. Por isso os \omterm{def:g10:universal-gravitation:satellite}{satélites} carregam relógios
atômicos --- e, nessa precisão, o ritmo de um relógio depende de com que velocidade ele
se move e de que altura ele ocupa, um efeito que encontramos no
\cref{ch:g12:special-relativity}.
\end{remark}

\section{Exercícios}

\begin{exercise}[$\star$]\label{exo:g12:oscillators-and-time:1}
(a) Calcule o \omterm{def:g12:oscillators-and-time:oscillator}{período} e a \omterm{def:g12:oscillators-and-time:oscillator}{frequência} de um \omterm{def:g12:oscillators-and-time:pendulum}{pêndulo simples} de \qty{1.00}{m}. (b) O
comprimento é quadruplicado: o que acontece com o \omterm{def:g12:oscillators-and-time:oscillator}{período}?
\end{exercise}

\begin{exercise}[$\star$]\label{exo:g12:oscillators-and-time:2}
Um coração em repouso bate $72$ vezes por minuto; a asa de uma libélula, a
\qty{30}{Hz}. Dê a \omterm{def:g12:oscillators-and-time:oscillator}{frequência} e o \omterm{def:g12:oscillators-and-time:oscillator}{período} do coração e o \omterm{def:g12:oscillators-and-time:oscillator}{período} da asa; que palavra do
capítulo nomeia a elevação máxima do peito a cada batida?
\end{exercise}

\begin{exercise}[$\star$]\label{exo:g12:oscillators-and-time:3}
Um carrinho de \qty{0.10}{kg} anda numa mola de \qty{25}{N/m}: calcule $T_0$ e $f_0$. O
que dobrar a massa faz com o \omterm{def:g12:oscillators-and-time:oscillator}{período}?
\end{exercise}

\begin{exercise}[$\star$]\label{exo:g12:oscillators-and-time:4}
Máximos sucessivos de um \omterm{def:g12:oscillators-and-time:oscillator}{oscilador} amortecido: $2.0$, $1.5$, $1.13$, \qty{0.85}{cm} em
$t = 0$, $2.0$, $4.0$, \qty{6.0}{s}. Leia o \omterm{def:g12:oscillators-and-time:damping}{pseudoperíodo}; mostre que a \omterm{def:g12:oscillators-and-time:oscillator}{amplitude} cai
por um fator fixo a cada \omterm{def:g12:oscillators-and-time:oscillator}{período}; preveja o próximo máximo.
\end{exercise}

\begin{exercise}[$\star$]\label{exo:g12:oscillators-and-time:5}
Uma astronauta leva um pêndulo e um \omterm{prop:g12:oscillators-and-time:spring-period}{oscilador massa--mola} para a Lua
($g = \qty{1.62}{m/s^2}$). Por que fator muda cada \omterm{def:g12:oscillators-and-time:oscillator}{período}? Explique.
\end{exercise}

\begin{exercise}[$\star\star$]\label{exo:g12:oscillators-and-time:6}
(a) Verifique que $\sqrt{L/g}$ tem dimensão de tempo. (b) Por que nenhuma fórmula para
$T_0$ construída com $m$, $L$ e $g$ pode conter a massa? (c) Esse argumento poderia
alguma vez produzir o $2\pi$?
\end{exercise}

\begin{exercise}[$\star\star$]\label{exo:g12:oscillators-and-time:7}
Verifique, derivando duas vezes, que $x(t) = A\cos(2\pi t/T_0 + \varphi)$ satisfaz
$\frac{d^2x}{dt^2} = -(k/m)\,x$ se e somente se $T_0 = 2\pi\sqrt{m/k}$.
\end{exercise}

\begin{exercise}[$\star\star$]\label{exo:g12:oscillators-and-time:8}
Uma mola de \omterm{def:g12:oscillators-and-time:spring}{constante elástica} \qty{40}{N/m} carrega um carrinho de \qty{0.20}{kg} com
\omterm{def:g12:oscillators-and-time:oscillator}{amplitude} de \qty{3.0}{cm}. Calcule a \omterm{def:g11:mechanical-energy:em}{energia mecânica} e a velocidade máxima. Onde a
velocidade é máxima? E nula?
\end{exercise}

\begin{exercise}[$\star\star$]\label{exo:g12:oscillators-and-time:9}
(a) Calcule o comprimento do pêndulo de segundos ($T_0 = \qty{2.00}{s}$). (b) O
comprimento dele cresce $1.0\%$: com $\sqrt{1+x} \approx 1 + x/2$, encontre a variação
relativa de $T_0$ e o erro diário.
\end{exercise}

\begin{exercise}[$\star\star$]\label{exo:g12:oscillators-and-time:10}
Um \omterm{def:g12:oscillators-and-time:oscillator}{oscilador} fracamente amortecido perde $5.0\%$ da \omterm{def:g10:energy-conservation:energy}{energia} dele a cada \omterm{def:g12:oscillators-and-time:oscillator}{período}. (a)
Qual é a queda de \omterm{def:g12:oscillators-and-time:oscillator}{amplitude} por \omterm{def:g12:oscillators-and-time:oscillator}{período}? (b) Quantos \omterm{def:g12:oscillators-and-time:oscillator}{períodos} para perder metade da
\omterm{def:g10:energy-conservation:energy}{energia}? (c) O \omterm{def:g12:oscillators-and-time:damping}{pseudoperíodo} está longe de $T_0$?
\end{exercise}

\begin{exercise}[$\star\star$]\label{exo:g12:oscillators-and-time:11}
Explique com a curva de \omterm{def:g12:oscillators-and-time:resonance}{ressonância}: (a) por que empurrar um balanço só funciona no
ritmo próprio dele; (b) por que os soldados quebram o passo em passarelas; (c) por que
uma máquina de lavar treme numa certa rotação enquanto acelera a centrifugação.
\end{exercise}

\begin{exercise}[$\star\star\star$]\label{exo:g12:oscillators-and-time:12}
Um relógio de pêndulo, exato ao nível do mar, muda-se para um observatório onde
$g = \qty{9.78}{m/s^2}$. (a) Ele adianta ou atrasa? Quantos segundos por dia? (b)
Quanto mais curto precisa ficar o pêndulo de \qty{0.994}{m} dele?
\end{exercise}

\begin{exercise}[$\star\star\star$]\label{exo:g12:oscillators-and-time:13}
Um carrinho de \qty{0.20}{kg} obedece a $x(t) = A\cos(2\pi t/T_0)$, com
$A = \qty{4.0}{cm}$ e $T_0 = \qty{0.63}{s}$. Calcule a velocidade máxima, a aceleração
máxima, a \omterm{def:g12:oscillators-and-time:spring}{constante elástica} $k$ e a \omterm{def:g11:mechanical-energy:em}{energia mecânica}.
\end{exercise}

\begin{exercise}[$\star\star\star$]\label{exo:g12:oscillators-and-time:14}
Para medir $g$: um pêndulo com $L = 99.5 \pm \qty{0.2}{cm}$ faz $20$ oscilações em
$40.1 \pm \qty{0.2}{s}$. Calcule $g$ e a \omterm{def:g10:orders-of-magnitude:uncertainty}{incerteza relativa} dele (some as de $L$ e de
$T_0^2$); o resultado concorda com \qty{9.81}{m/s^2}?
\end{exercise}

\begin{exercise}[$\star\star\star$]\label{exo:g12:oscillators-and-time:15}
Compare marcadores de tempo pelo desvio relativo (erro dividido pelo tempo decorrido):
um relógio de pêndulo que atrasa \qty{10}{s} por dia, um relógio de quartzo que erra
\qty{0.5}{s} por mês e um relógio de césio que erra \qty{1}{s} em cem milhões de anos:
calcule as três razões. Um relógio de GPS errado em \qty{1}{\micro s}: qual é o erro de
posição? De que família o GPS precisa?
\end{exercise}

\section{Problema: a bancada do relojoeiro}

\begin{problem}[{Problema de fim de semana --- a bancada do relojoeiro: um relógio de
parede trazido de volta à hora certa --- cinquenta balanços cronometrados, uma haste de
latão que estica no verão, o empurrãozinho sussurrado do escapamento, um desafiante de
quartzo --- e o veredito dado em segundos por mês}]
\label{pb:g12:oscillators-and-time:1}
Um relógio de parede de pêndulo chega à bancada ``andando errado''. Trate a haste de
latão e o corpo pesado dele como um \omterm{def:g12:oscillators-and-time:pendulum}{pêndulo simples} de comprimento efetivo $L$, com
massa $m = \qty{1.2}{kg}$; a engrenagem conta cada oscilação completa como exatamente
\qty{2.000}{s}.

\medskip
\noindent\textbf{Parte I --- Tomar o pulso do relógio.}
\begin{enumerate}
  \item Por que cronometrar $50$ oscilações, e não uma? Estime o erro que sobra em
        $T_0$ se o seu tempo de reação for de \qty{0.2}{s}.
  \item O cronômetro marca \qty{99.4}{s} para $50$ oscilações: qual é o \omterm{def:g12:oscillators-and-time:oscillator}{período}?
  \item Deduza o comprimento efetivo atual do pêndulo.
  \item Que comprimento daria exatamente \qty{2.000}{s}? A porca de regulagem deve
        subir ou descer o corpo, e de quanto?
  \item Sem ajuste, o relógio adianta ou atrasa, e quantos segundos por dia?
  \item Sensibilidade a $g$: qual seria o \omterm{def:g12:oscillators-and-time:oscillator}{período} na Lua ($g = \qty{1.62}{m/s^2}$)? E o
        erro diário num observatório onde $g = \qty{9.80}{m/s^2}$?
\end{enumerate}

\medskip
\noindent\textbf{Parte II --- A moleza do verão.} O latão se dilata: um aquecimento
$\Delta\theta$ estica a haste até $L' = L\,(1 + \lambda\, \Delta\theta)$, com
$\lambda = \num{1.9e-5}$ por \omterm{def:g10:light-spectra:kelvin}{kelvin}. O relógio é ajustado a \qty{20}{\celsius}.
\begin{enumerate}[resume]
  \item Calcule $\Delta L$ para uma sala de verão a \qty{25}{\celsius}.
  \item Usando $\sqrt{1+x} \approx 1 + x/2$, mostre que
        $\Delta T_0 / T_0 = \tfrac12 \lambda\, \Delta\theta$.
  \item Calcule $\Delta T_0/T_0$ e o erro diário de verão: adianta ou atrasa?
  \item As mesmas perguntas para um corredor de inverno a \qty{10}{\celsius}.
  \item Os antigos relógios de precisão usavam pêndulos de ``grelha'', que misturavam
        hastes de latão com hastes de aço (que se dilatam cerca da metade) em sentidos
        alternados. Explique como isso pode segurar o ritmo.
\end{enumerate}

\medskip
\noindent\textbf{Parte III --- O empurrão sussurrado.} Com o escapamento desengatado, o
balanço, iniciado com \omterm{def:g12:oscillators-and-time:oscillator}{amplitude} $\theta_m = 2.00^\circ$, reduz a \omterm{def:g12:oscillators-and-time:oscillator}{amplitude} à metade em
\qty{280}{s}; em serviço, um empurrãozinho do escapamento a cada \omterm{def:g12:oscillators-and-time:oscillator}{período} mantém
$\theta_m$ constante.
\begin{enumerate}[resume]
  \item Quantos \omterm{def:g12:oscillators-and-time:oscillator}{períodos} leva essa redução à metade?
  \item No extremo o corpo subiu $h = L(1 - \cos\theta_m) \approx L\,\theta_m^2/2$ (em
        radianos): deduza $E_m = \tfrac12 m g L\,\theta_m^2$ e calcule-a
        (\cref{ch:g11:mechanical-energy}).
  \item A \omterm{def:g10:energy-conservation:energy}{energia} é proporcional ao quadrado da \omterm{def:g12:oscillators-and-time:oscillator}{amplitude}: que fração de $E_m$ se perde
        por \omterm{def:g12:oscillators-and-time:oscillator}{período}?
  \item Deduza a \omterm{def:g10:energy-conservation:energy}{energia} de cada empurrão e a \omterm{def:g11:work-of-force:power}{potência} média necessária.
  \item O \omterm{def:g10:universal-gravitation:weight}{peso} motor $M = \qty{2.5}{kg}$ desce $h = \qty{1.0}{m}$ por semana: isso é
        \omterm{def:g10:energy-conservation:energy}{energia} suficiente, e com que \omterm{def:g10:energy-conservation:efficiency}{rendimento}?
\end{enumerate}

\medskip
\noindent\textbf{Parte IV --- O desafiante de quartzo.} Um movimento de quartzo oscila a
$f = \qty{32768}{Hz}$ e erra cerca de \qty{0.3}{s} por mês; o pêndulo restaurado ainda
erra \qty{0.5}{s} por dia.
\begin{enumerate}[resume]
  \item Calcule o \omterm{def:g12:oscillators-and-time:oscillator}{período} do quartzo. Mostre que $\num{32768} = 2^{15}$: por que uma
        potência de dois é exatamente o que o circuito de um relógio quer?
  \item Calcule os desvios relativos do quartzo e do pêndulo restaurado, e cada um em
        segundos por mês. Quem vence, e por que fator?
  \item O padrão de césio --- \num{9192631770} \omterm{def:g12:oscillators-and-time:oscillator}{períodos} fazem um segundo --- erra
        \qty{1}{s} em cem milhões de anos: qual é o desvio relativo dele? Quantas \omterm{def:g10:orders-of-magnitude:oom}{ordens
        de grandeza} abaixo do quartzo, e por que o GPS o exige?
  \item O veredito, em uma frase com números: o que o relógio restaurado ganha na
        parede, e a quem pertence agora o segundo?
\end{enumerate}
\end{problem}
